% ID: SEL-D0063-OA-E02
% ESTADO: cotejada. Segunda pasada visual y borrador compilado revisados.
% Fuente: Selectividad/sel_2010_matematicas_aplicadas/mat_apli/mat_apli_junio/mat_apli_junio_rb/EXAMEN-ANDALUCÍA-6.pdf
% SHA-256: 0de2f7d077a3b5d5c0dbca886cdb363082356411843369c3bce54ff7cda08abc
% Página: 1; opción A; ejercicio original 2.

Sea la función $f(x)=2x^2+ax+b$.
\begin{enumerate}[label=\alph*)]
\item \textbf{(1.25 puntos)} Determine los valores de $a$ y $b$ sabiendo que su gráfica pasa por el punto $(1,3)$ y alcanza un extremo local en el punto de abscisa $x=-2$.
\item \textbf{(1.25 puntos)} Tomando $a=8$ y $b=-10$ deduzca la curvatura de su gráfica, el valor mínimo que alcanza la función y los valores donde la función se anula.
\end{enumerate}
